\section{Notation, Definition, and Representation} \label{Appendix A: Notation, Definition}
\subsection{Notation}
\paragraph{Acronym}~\\
TV: Total Variation \\
BTV: Bounded Total Variation \\
SVN: Sectional Variational Norm \\
BSVN: Bounded Sectional Variational Norm \\
HAL: Highly Adaptive Lasso \\
TMLE: Target Maximum Likelihood Estimator (or Minimum Loss)\\
HAL-TMLE: TMLE that uses HAL for initial estimation \\
HAL-MLE: Highly Adaptive Lasso Maximum Likelihood Estimator \\
KDE: Kernel Density Estimation \\
LAS: Local Adaptive Splines \\
TF: Trend Filtering \\
TFPP: Trend Filtering with Penalty on the Parametric Part \\
LSDE: Logspline Density Estimation \\
PLSDE: TV-penalized Logspline Density Estimation \\
KM: Kaplan-Meier Estimator \\
CV: Cross Validation \\
càdlàg function: multivariate real-valued function on (say) unit cube $[0, 1]^d$ that is right-continuous with left-hand limits \\
CDF: generalized cumulative distribution function.\\
DGP: Data-Generating Process. \\
LLFP: local least favorable path \\
ULFP: universal least favorable path \\
FISTA: Fast Iterative Shrinkage-Thresholding Algorithm \\
L-BFGS: Limited-memory Broyden-Fletcher-Goldfarb-Shanno \\
AdaGrad: Adaptive Gradient \\

\paragraph{Mathematical Notation}~\\
$\sim$: We say $X \sim P_0$ if the random variable $X$ has distribution $P_0$. \\
$\lesssim$: We say $a \lesssim b$ if $a \leq Cb$ for some constant $C$.\\
$\asymp$: We say $a \asymp b$ if $cb \leq a \leq Cb$ for some constant $c$ and $C$. \\
$\asymp^+$: We say $a \asymp^+ b$ if $cb \leq a \leq 
Cb(\log n)^p$ for some constant $c$, $C$, and finite $p$. Note that $p$ can be negative. \\
$\equiv$: Definition.
$\mathrm{TV}(\cdot)$: Total variation functional as defined in Section 9.1 of \citet{tibshirani2022divided}. When the function $f$ is a càdlàg function, $\mathrm{TV}(f) = \int_0^1 |df(u)|$. This is also referred to as the total variation norm (a seminorm), denoted as $|f|_v$. \\
$X$: Euclidean valued observed random variable. \\
$P_0$: True data-generating distribution of X or the truth expectation operator in empirical process. \\
$P_n$: Empirical data-generating distribution of X or the empirical expectation operator in empirical process. \\
$P_{f_n}$: Expectation operator with measure $Q_n$. \\
$P_{f_0}$: Expectation operator with measure $f_0$, which is the truth $P_0$. \\
$(\mathbb{G}_n(f) : f \in \mathcal{F})$: denotes empirical process $\mathbb{G}_n = n^{1/2}(P_n - P) \in \ell^{\infty}(\mathcal{F})$ indexed by a class of functions $\mathcal{F}$.\\
$\mathcal{M}$: Statistical model, the set of possible probability distributions for $P_0$.\\
$O^+(r(n))$: We use notation $O^+(r(n))$ if the term is $O(r(n)(\log n)^p)$ for some finite $p$, and the $+$ notation is also used in $O_p^+$. In words, it is described as "up till a $\log n$ factor".\\
$d_0(f, f_0)$: loss-based dissimilarity between $f$ and $f_0$ defined by $d_0(f, f_0) = P_0L(f) - P_0L(f_0)$ for some general loss $L$.\\
$P_0 \in \mathcal{M}$: $P_0$ is known to be an element of the statistical model $\mathcal{M}$.\\
$\mathrm{d}f/\mathrm{d}\mu$: Radon-Nikodym derivative of measure implied by cadlag function $f$ w.r.t. Lebesgue measure $\mu$.\\
$p = \mathrm{d}P/\mathrm{d}\mu$: Density of $P$ w.r.t. dominating measure $\mu$\\
$p_0$: True density of data-generating distribution $P_0$ w.r.t. appropriate dominating measure $\mu$\\
$p_f$: The density $p$ indexed by a function $f$ through a link function.\\
$\| f \|_v$: variation norm of cadlag function $f$.\\
$f^{(k)}$: k-th order Lebesgue-Radon-Nikodym derivative of cadlag function $f$. \\
$N(\varepsilon, \mathcal{F}, \| \cdot \|)$: the covering number of class $\mathcal{F}$ w.r.t. some metric defined as the minimal number of balls with radius $\varepsilon$ needed to cover $\mathcal{F}$.\\
$J(\delta, \mathcal{F}, \| \cdot \|)$: the entropy integral for class $\mathcal{F}$ w.r.t. some metric defined as $\int_0^{\delta} \sqrt{\log N(\varepsilon, \mathcal{F}, \| \cdot \|)} \mathrm{d}\varepsilon$. $J_{\infty}(\delta, \mathcal{F})$ is used when the norm is the supremum norm. $J_{2}(\delta, \mathcal{F})$ is used when the norm is the $L^2$ norm.\\
$\Pi$: projection operator.\\
$S_f$: score operator at $f$.\\
$D^*(P)$: canonical gradient (efficient influence curve) of a pathwise differentiable parameter at $P$.\\

\paragraph{Definition}~\\
$f_n$: we use subscript n for estimations, like $\beta_n$ or $C_n$. Here in this paper, $f_n$ means the HAL-MLE.\\
$f_{n, \beta_n(M)}$: the HAL-MLE with a user-specified M. So it is parameterized by the coefficient $\beta_n$ under the $L_1$ constraint less than M. $f_{n, \beta_n(M)}$ is the MLE over $D^{(k)}_M(\mathcal{R}_n)$\\
$M_{cv}$: the hyperparameter M chosen by CV.\\
$f_n^{rel}$: Relaxed HAL-MLE. Relax the HAL-MLE $f_{n, \beta_n(M)}$ by loosening the $L_1$ constraint with a larger constant $C > M$, such that it is a MLE over $D^{(k)}_C(\mathcal{R}_n)$. Notice that we use the word "fully relaxed" when we drop the $L_1$ constraint completely \\
Undersmooth: The behavior that we intentionally include more basis to reduce the bias, which will increase variance. Practically, we can choose $M' > M_{cv}$ to realize this. It represents the trade off between bias and variance. \\
$f_{n,0}$: The oracle HAL-MLE over a sample of size n. Such that $f_{n,0}=\arg\min_{f\in D^{(k)}_M({\cal R}_{n,0} )}-P_0 \log p_{f}$. \\
$\tilde f_n$: The projection of $f_n$ to some other linear subspace. Typically, we used $\tilde f_n$ to denote the projection of $f_n$ on $ D^{(k)}_M({\cal R}_{n,0})$ in Appendix~\ref{appendix_D: proof of asymptotic normality}. \\
$f_n^k$: the k-th step TMLE update, shorthand notation of $f_{n, \varepsilon_n^k}$. $f_n$ is also $f_n^0$. \\
$f^{(k)}$: the k-th order derivative of f. \\
$C_n$: The normalizing constant of $\tilde f_n$, also denoted as $C(\tilde f_n)$, or $C(\tilde \beta_n)$. \\
$C_{n,0}$: The oracle normalizing constant, also denoted as $C(f_{n,0})$, or $C(\beta_{n,0})$. \\
$\| f \|_v^*$: sectional variation norm $\| f \|_v^* = \int_{[0,1]^d} | df(u) |$.\\
$D^{(0)}_U([0, 1]^d)$: Space of real valued cadlag functions on $[0, 1]^d$ with BSVN for some $U < \infty$. We use $D^{(0)}_U([0, 1])$ for our univariate case.\\
$\| f \|_{v,k}^*$: $k$-th order sectional variation norm of f.\\
$D^{(k)}_U([0, 1])$: functions $f$ in $D^{(0)}_U([0, 1])$ for which all k-th order Lebesgue-Radon-Nikodym derivatives exist and $\| f \|_{v,k}^* < U < \infty$\\
$\mathcal{R}_0$: complete infinite index set of basis.\\
$\mathcal{R}_n$: a finite data-adaptive index set of basis. The active set after Lasso selection. For the nonparametric part of the basis system, we index them each by a tuple $(k, u)$ for $k$-th order splines and any knot point $u \in (0,1]$. For the parametric part of the basis system, we index them each by a tuple $(i, 0)$, for $i \leq k$ and starting point 0. As such, we could define the total index set as a union of all these. \\
$D^{(k)}(\mathcal{R}_n)$: for a given subset $\mathcal{R}_n \subset \mathcal{R}_0$, this is defined as the space of functions that are in closure of linear span of $\{\phi_{k,u} : (k, u) \in \mathcal{R}_n\}$. \\
$D^{(k)}_M(\mathcal{R}_n)$: the subset of $D^{(k)}(\mathcal{R}_n)$ with a user-specified $M$ as the $L_1$ norm constraint of the basis.\\
$D^{(k)}_M(\mathcal{R}_{n,0})$: an independent oracle space with a user-specified $M$ as the $L_1$ norm constraint of the basis. A technique used for proving pointwise asymptotic normality.\\
$J_n$: the effective dimension of $D^{(k)}_M(\mathcal{R}_n)$. The number of active basis. \\
$J_{n,0}$: the effective dimension of $D^{(k)}_M(\mathcal{R}_{n,0})$. The number of active basis. \\
$\phi_j^*$: some orthonormal basis function on $D^{(k)}_M(\mathcal{R}_{n,0})$.\\
$\bar \phi_{n,x}$: a special basis function at point $x$ defined as $\bar \phi_{n,x} = J_{n,0}^{-1/2} \sum_{j \in \mathcal{R}_{n,0}} \phi_j^*(x) \phi_j^*$. Used interchangably with $\bar \phi$.\\
$\tilde \phi_j$: the projection of $\phi_j^*$ onto some linear span, like $D^{(k)}_M(\mathcal{R}_n)$. In general, we use the tilde sign above an element to denote its projection on some other spaces, e.g. $\tilde \phi_{n,x}$, $\tilde f_n$.\\




\subsection{Representation Theorem} \label{App:representation thm}
This is the proof for Proposition \ref{thm:HAL_representation}

\begin{proof}
\textit{Base case:} For $k = 0$ (with $u_0 = u$),
\[
  f(x)
  = f(0) + \int_{(0, x]} \mathrm{d}\,f(u)
  = f(0) + \int_{(0, 1]} I\{u \le x\}\,\mathrm{d}\,f(u)
  = f(0) + \int_0^1 (x - u)_+^0 \,\mathrm{d}\,f(u).
\]

\textit{Induction:} If, for a $k$‑times differentiable function $f$, the representation holds true, we now show it for a $(k+1)$‑times differentiable $f$.

Since $f^{(k)}$ is differentiable,
\[
  \mathrm{d}\,f^{(k)}(u_k) = f^{(k+1)}(u_k)\,\mathrm{d}u_k,
  \quad f^{(k+1)}(u_k)\in D^{(0)}_U([0,1]).
\]
Starting from the inductive hypothesis,
\begin{align*}\label{HAL_representation_formula}
  f(x)
  &= \sum_{i = 0}^k \frac{1}{i!}f^{(i)}(0)\,x^{i}
    + \int_0^1 \frac{1}{k!}(x - u_k)_+^k\,\mathrm{d}\,f^{(k)}(u_k) \\
  &= \sum_{i = 0}^k \frac{1}{i!}f^{(i)}(0)\,x^{i}
    + \int_0^1 \frac{1}{k!}(x - u_k)_+^k\,f^{(k+1)}(u_k)\,\mathrm{d}u_k \\
  &= \sum_{i = 0}^k \frac{1}{i!}f^{(i)}(0)\,x^{i}
    + \int_0^1 \frac{1}{k!}(x - u_k)_+^k \\
  &\quad \times \Bigl(f^{(k+1)}(0)
    + \int_0^{u_k} \mathrm{d}\,f^{(k+1)}(u_{k+1})\Bigr)\,\mathrm{d}u_k \\
  &= \sum_{i = 0}^k \frac{1}{i!}f^{(i)}(0)\,x^{i}
    + \int_0^1 \frac{1}{k!}(x - u_k)_+^k\,f^{(k+1)}(0)\,\mathrm{d}u_k \\
  &\quad + \int_0^1 \frac{1}{k!}\int_0^{u_k}(x - u_k)_+^k\,\mathrm{d}\,f^{(k+1)}(u_{k+1})\,\mathrm{d}u_k \\
  &= \sum_{i = 0}^k \frac{1}{i!}f^{(i)}(0)\,x^{i}
    + f^{(k+1)}(0)\int_0^x \frac{1}{k!}(x - u_k)^k\,\mathrm{d}u_k \\
  &\quad + \int_0^1 \frac{1}{k!}\int_0^{u_k}(x - u_k)_+^k\,\mathrm{d}\,f^{(k+1)}(u_{k+1})\,\mathrm{d}u_k \\
  &= \sum_{i = 0}^{k+1} \frac{1}{i!}f^{(i)}(0)\,x^{i}
    + \int_0^1 \frac{1}{k!}\int_0^{u_k}(x - u_k)_+^k\,\mathrm{d}\,f^{(k+1)}(u_{k+1})\,\mathrm{d}u_k \\
  &= \sum_{i = 0}^{k+1} \frac{1}{i!}f^{(i)}(0)\,x^{i} \\
  &\quad + \int_0^1 \frac{1}{k!}\int_0^1(x - u_k)^k I\{u_{k+1}\le u_k\le x\}\,\mathrm{d}\,f^{(k+1)}(u_{k+1})\,\mathrm{d}u_k \\
  &= \sum_{i = 0}^{k+1} \frac{1}{i!}f^{(i)}(0)\,x^{i} \\
  &\quad + \int_0^1 \frac{1}{k!}\int_0^1(x - u_k)^k I\{u_{k+1}\le u_k\le x\}\,\mathrm{d}u_k\,\mathrm{d}\,f^{(k+1)}(u_{k+1}) \\
  &\qquad\text{(by Fubini's theorem)}\\
  &= \sum_{i = 0}^{k+1} \frac{1}{i!}f^{(i)}(0)\,x^{i}
    + \int_0^1\frac{1}{k!}\Bigl(\int_{u_{k+1}}^{x}(x - u_k)^k\,\mathrm{d}u_k\Bigr)\,\mathrm{d}\,f^{(k+1)}(u_{k+1})\\
  &= \sum_{i = 0}^{k+1} \frac{1}{i!}f^{(i)}(0)\,x^{i}
    + \int_0^1\frac{1}{(k+1)!}(x - u_{k+1})_+^{k+1}\,\mathrm{d}\,f^{(k+1)}(u_{k+1}).
\end{align*}

Then, by induction, the representation holds for all \(k\).
\end{proof}